\chapter{The original obstacle sequence on the actual projective family}
\label{chap:c0n-graded-base-ideals}

This chapter retains the original projective subvariety $X\subset\mathbf P^M_k$,
its original homogeneous prime ideal, the original obstacle $\phi$, and the
original ideals $J_m(\phi)$. The affine cone may be singular: multiplier ideals
are only formed on the actual smooth projective family or its smooth charts.
No desired gradedness, nonvanishing, generation, or cone comparison is an input.

\begin{lemma}[Whole homogeneous ideal-sheaf arithmetic]
\label{lem:c0n-homogeneous-sheaf-arithmetic}
\lean{EnvelopeProperty.Char0.EnvelopeGradedBaseIdealsCore.quotientIdealSheaf_coordinateIdeal, EnvelopeProperty.Char0.EnvelopeGradedBaseIdealsCore.quotientIdealSheaf_mono, EnvelopeProperty.Char0.EnvelopeGradedBaseIdealsCore.quotientIdealSheaf_mul, EnvelopeProperty.Char0.EnvelopeGradedBaseIdealsCore.quotientIdealSheaf_top}
\leanok
For the ordinary projective scheme over any commutative ring, sheafification of
homogeneous ideals preserves inclusion, products, and the unit ideal. Every
coordinate ideal is the actual dehomogenization image under the original
native coordinate-ring isomorphism.
\end{lemma}
\begin{proof}
The native homogeneous quotient closed immersion has kernel ideal whose
pullback under the coordinate-ring isomorphism is exactly the original
homogeneous ideal mapped by dehomogenization. Applying the forward ring
isomorphism gives the entire coordinate ideal, with no radical or saturation
replacement. Ring-ideal image preserves inclusions, products and the unit
ideal. The standard coordinate opens cover the whole native projective scheme;
ideal-sheaf detection on this actual cover proves the three whole-sheaf claims.
\end{proof}

\begin{lemma}[The same obstacle polynomial with parameter in the coefficient ring]
\label{lem:c0n-obstacle-parameter-swap}
\lean{EnvelopeProperty.ProjSubvariety.MultiplierBaseIdeals.parameterEquiv, EnvelopeProperty.ProjSubvariety.MultiplierBaseIdeals.parameterEquiv_C, EnvelopeProperty.ProjSubvariety.MultiplierBaseIdeals.parameterEquiv_X, EnvelopeProperty.ProjSubvariety.MultiplierBaseIdeals.parameterEquiv_monomial, EnvelopeProperty.ProjSubvariety.MultiplierBaseIdeals.parameterObstacleIdeal, EnvelopeProperty.ProjSubvariety.MultiplierBaseIdeals.parameterObstacleIdeal_isHomogeneous, EnvelopeProperty.ProjSubvariety.MultiplierBaseIdeals.homogeneousObstacleIdeal, EnvelopeProperty.ProjSubvariety.MultiplierBaseIdeals.homogeneousObstacleIdeal_mul_le, EnvelopeProperty.ProjSubvariety.MultiplierBaseIdeals.homogeneousObstacleIdeal_zero}
\leanok
There is a canonical $k$-algebra isomorphism
$k[x_0,\ldots,x_M][t]\simeq k[t][x_0,\ldots,x_M]$ fixing every $x_i$ and $t$.
The image of $J_m(\phi)$ is homogeneous in the original $x$-grading. Its
bihomogeneous generators still have $x$-degree $md$ and their original
parameter exponents.
\end{lemma}
\begin{proof}
Compose the inverse of the native option-variable polynomial equivalence with
its coefficient-polynomial version. The resulting map sends a coefficient
form $s$ to its coefficient extension and sends the original parameter to the
constant polynomial $t$. The original obstacle ideal is generated by all
monomials $s t^l$ belonging to its original obstacle space. Their images are
$t^l s$, homogeneous of degree $md$. An ideal generated by homogeneous elements
is homogeneous. Thus this construction retains the entire original ideal,
not a selected subfamily or an assumed presentation.
\end{proof}

\begin{proposition}[Actual original projective-family ideal sequence]
\label{prop:c0n-original-obstacle-scheme-sequence}
\lean{EnvelopeProperty.ProjSubvariety.MultiplierBaseIdeals.originalHomogeneousIdeal, EnvelopeProperty.ProjSubvariety.MultiplierBaseIdeals.familyHomogeneousIdeal, EnvelopeProperty.ProjSubvariety.MultiplierBaseIdeals.familyScheme, EnvelopeProperty.ProjSubvariety.MultiplierBaseIdeals.familyInclusion, EnvelopeProperty.ProjSubvariety.MultiplierBaseIdeals.familyInclusion_isClosedImmersion, EnvelopeProperty.ProjSubvariety.MultiplierBaseIdeals.familyProjection, EnvelopeProperty.ProjSubvariety.MultiplierBaseIdeals.familyProjection_isProjectiveFactorization, EnvelopeProperty.ProjSubvariety.MultiplierBaseIdeals.ambientSchemeIdeal, EnvelopeProperty.ProjSubvariety.MultiplierBaseIdeals.schemeIdeal, EnvelopeProperty.ProjSubvariety.MultiplierBaseIdeals.schemeIdeal_mul_le, EnvelopeProperty.ProjSubvariety.MultiplierBaseIdeals.schemeIdeal_zero, EnvelopeProperty.ProjSubvariety.MultiplierBaseIdeals.ambientSchemeIdeal_coordinate}
\leanok
\uses{lem:c0n-homogeneous-sheaf-arithmetic,lem:c0n-obstacle-parameter-swap}
Let $\mathcal X=X\times\mathbf A^1_k$ be the original homogeneous quotient over
$k[t]$, with its actual closed inclusion in $\mathbf P^M_{k[t]}$. Pulling back
the sheaf ideals of the parameter-swapped $J_m(\phi)$ defines whole ideals
$\mathfrak a_m$ on this same $\mathcal X$. They satisfy
$\mathfrak a_0=\mathcal O_{\mathcal X}$ and
$\mathfrak a_m\mathfrak a_r\subseteq\mathfrak a_{m+r}$.
The ambient ideal on every original coordinate chart is the full original
homogeneous dehomogenization ideal; the family ideal is its actual geometric
comap along the original closed inclusion.
\end{proposition}
\begin{proof}
Use the original homogeneous prime ideal extended along $k\to k[t]$ for the
quotient; no model or comparison is supplied. Apply the preceding sheaf
arithmetic to the actual original obstacle product inclusion
$J_mJ_r\subseteq J_{m+r}$ and $J_0=(1)$, then use the native geometric ideal
comap, which preserves products and the unit ideal. The ambient coordinate formula
is the same original quotient-kernel calculation. Computing its actual
closed-family comap on native quotient charts remains part of the next
nonvanishing and generation producer.
\end{proof}

\begin{proposition}[Nonzero original ideals and actual line generation]
\label{prop:c0n-original-obstacle-nonzero-generation}
\uses{prop:c0n-original-obstacle-scheme-sequence}
If $\phi\ge C$, the original sequence has nonzero generic stalks in every
index. Its same whole ideal $\mathfrak a_m$ tensored by the actual
$\mathcal O_{\mathcal X}(md)$ is generated by the original obstacle monomial
sections; on a Noetherian projective family a finite original subfamily suffices.
Consequently any required positive Veronese is selected internally.
\end{proposition}
\begin{proof}
Set $N=\lceil-mC\rceil_+$. The original obstacle space contains
$x_i^{md}t^N$ for every coordinate $i$. On a nonempty $x_i$ chart its normalized
coefficient is exactly $t^N$, which is nonzero in the original domain
$\mathcal O_X(U_i)[t]$. Injectivity into the original generic field proves
nonzero generic stalk, with no nonvanishing certificate. Every generator
$s t^l$ is the actual global section of $\mathcal O(md)$ defined by its same
homogeneous form and parameter coefficient. The canonical dual evaluation
has precisely the original dehomogenization ideal on every coordinate chart.
Its true image factor, tensored with the original line and contracted by the
canonical evaluation isomorphism, is an epimorphism onto
$\mathfrak a_m\otimes\mathcal O(md)$. The free module's actual global generators
therefore generate this whole original module. Finite presentation over the
Noetherian projective family selects finitely many original sections. A
Veronese uses these very indices and the native tensor-power multiplication
isomorphism; it does not impose generation of a new sequence.
\end{proof}

\begin{proposition}[The actual multiplier cone ideal]
\label{prop:c0n-original-multiplier-cone-ideal}
\uses{prop:c0n-original-obstacle-nonzero-generation}
Define $B_m$ by taking the preimages of the actual affine multiplier ideals of
$\mathfrak a_\bullet$ on all original smooth charts and intersecting in the
original cone ring. This same ideal contains $J_m$, satisfies
$B_{ml}\subseteq(B_m^l)^{\mathrm{sat}}$, and has homogeneous generators of
$x$-degree at most $md+c_0$ after the internally produced uniform Nadel twist.
\end{proposition}
\begin{proof}
The whole scheme multiplier ideal restricts to the full affine multiplier
ideal under the original generic-field section equivalence. The original
ideal containment and genuine multiplier subadditivity imply the first two
claims on each original chart. The chart cover detects the corresponding
whole projective ideals, hence their cone relations modulo irrelevant
saturation. Uniform generation is obtained by applying the explicit ordinary
Nadel input and the actual Castelnuovo--Mumford argument to the preceding
original finite generators. Global sections of the resulting original line
twist have the same homogeneous coordinates; taking their images supplies
$G_m$ with $\langle G_m\rangle\subseteq B_m\subseteq
\langle G_m\rangle^{\mathrm{sat}}$. This step still requires the actual
scheme-to-cone coordinate comparison and is not an assumed generator bound.
\end{proof}

\begin{lemma}[Integrality of the same original parameter family]
\label{lem:c0n-original-family-integral}
\lean{EnvelopeProperty.ProjSubvariety.MultiplierBaseIdeals.familyIdeal_eq_parameterMap, EnvelopeProperty.ProjSubvariety.MultiplierBaseIdeals.familyIdeal_isPrime, EnvelopeProperty.ProjSubvariety.MultiplierBaseIdeals.coordinate_notMem_familyIdeal, EnvelopeProperty.ProjSubvariety.MultiplierBaseIdeals.family_irrelevant_ne_bot, EnvelopeProperty.ProjSubvariety.MultiplierBaseIdeals.familyScheme_isIntegral}
\leanok
\uses{lem:c0n-obstacle-parameter-swap,prop:c0n-original-obstacle-scheme-sequence}
The actual defining ideal of the original parameter-family quotient is
$E(I\cdot k[x][t])$, where $E$ is the original parameter exchange. It is prime;
a coordinate originally outside $I$ remains outside it, and the same actual
quotient Proj is integral.
\end{lemma}
\begin{proof}
The coefficient extension and the composite $k[x]\to k[x][t]\xrightarrow E
k[t][x]$ agree on every original polynomial. Ring-ideal image composition
therefore gives the full defining-ideal equality. A prime ideal remains prime
under ordinary polynomial extension; the actual ring isomorphism $E$ preserves
primality. If an original coordinate entered the family ideal, applying
$E^{-1}$ and taking the zeroth parameter coefficient would put that same
coordinate in $I$, a contradiction. Its nonzero degree-one class lies in the
irrelevant ideal, so that ideal is nonzero. The native integral-Proj theorem
for a domain with nonzero irrelevant ideal proves integrality of this same
family, including nonemptiness. No integral model is assumed.
\end{proof}

\begin{lemma}[Finitely many of the original homogeneous obstacle forms]
\label{lem:c0n-original-finite-obstacle-forms}
\lean{EnvelopeProperty.ProjSubvariety.MultiplierBaseIdeals.originalGeneratorSet, EnvelopeProperty.ProjSubvariety.MultiplierBaseIdeals.parameterObstacleIdeal_eq_span_originalGeneratorSet, EnvelopeProperty.ProjSubvariety.MultiplierBaseIdeals.originalGeneratorSet_homogeneous, EnvelopeProperty.ProjSubvariety.MultiplierBaseIdeals.exists_finite_original_homogeneous_generators}
\leanok
\uses{lem:c0n-obstacle-parameter-swap}
For every original index $m$, finitely many of the same original obstacle
monomials generate the full parameter-swapped ideal. Every selected form has
exactly the original $x$-degree $md$.
\end{lemma}
\begin{proof}
The original obstacle monomial decomposition gives the whole ideal as the span
of the images of all original $s t^l$ in the obstacle space. Each image is
$t^l s$ and retains the degree $md$. The ring $k[t][x_0,\ldots,x_M]$ is
Noetherian. Its ideal span is therefore finitely generated, and the native
finite-span-subset theorem selects a finite subset of this very generating
set. The resulting finite set is not an arbitrary unrelated presentation;
its subset proof retains membership in the original obstacle family.
\end{proof}

\begin{proposition}[The actual original generic stalks are unit ideals]
\label{prop:c0n-original-obstacle-generic-stalk}
\lean{EnvelopeProperty.ProjSubvariety.MultiplierBaseIdeals.parameterEquiv_mem_familyIdeal_iff, EnvelopeProperty.ProjSubvariety.MultiplierBaseIdeals.family_genericPoint_eq_zeroPoint, EnvelopeProperty.ProjSubvariety.MultiplierBaseIdeals.familyInclusion_zeroPoint_mem_iff, EnvelopeProperty.ProjSubvariety.MultiplierBaseIdeals.schemeIdeal_generic_stalk_eq_top, EnvelopeProperty.ProjSubvariety.MultiplierBaseIdeals.schemeIdeal_generic_stalk_ne_bot}
\leanok
\uses{lem:c0n-original-family-integral,prop:c0n-original-obstacle-scheme-sequence}
If the original obstacle has lower bound $C$, then every whole original
$\mathfrak a_m$ has unit, hence nonzero, generic stalk on the same actual
integral parameter family, including $m=0$. No Veronese or generic
nonvanishing premise is needed for this sequence.
\end{proposition}
\begin{proof}
The native generic point of the quotient Proj is the zero homogeneous prime.
The original closed inclusion sends this point to the prime of the original
family equations. A polynomial $g\in k[x][t]$ has image in that defining ideal
if and only if every original parameter coefficient of $g$ lies in $I$; this
is the genuine polynomial-extension ideal criterion transported by the actual
parameter isomorphism. Choose an original relevant coordinate $x_j\notin I$,
put $N=\lceil-mC\rceil_+$, and take the original obstacle generator
$g=x_j^{md}t^N$. It belongs to the original obstacle space. Its coefficient at
$N$ is $x_j^{md}\notin I$, since $I$ is prime. Thus its image is outside the
family defining prime. The original whole quotient-kernel ideal sheaf has
support equal to the projective zero locus of the entire original ideal.
The actual geometric comap pulls back that support. The preceding generator
shows that the same family generic point is outside it, so the true stalk
ideal is the unit ideal. Integrality makes the stalk ring nontrivial, hence
this unit ideal is not zero.
\end{proof}

\begin{lemma}[Canonical cone representative of a whole original scheme ideal]
\label{lem:c0n-canonical-scheme-cone-representative}
\lean{EnvelopeProperty.ProjSubvariety.MultiplierBaseIdeals.nativeRepresentative_mono, EnvelopeProperty.ProjSubvariety.MultiplierBaseIdeals.homogeneousIdeal_le_nativeRepresentative, EnvelopeProperty.ProjSubvariety.MultiplierBaseIdeals.parameterConeQuotient, EnvelopeProperty.ProjSubvariety.MultiplierBaseIdeals.parameterConeQuotient_parameterEquiv, EnvelopeProperty.ProjSubvariety.MultiplierBaseIdeals.homogeneousConeRepresentative, EnvelopeProperty.ProjSubvariety.MultiplierBaseIdeals.homogeneousConeRepresentative_comap, EnvelopeProperty.ProjSubvariety.MultiplierBaseIdeals.coneIdeal, EnvelopeProperty.ProjSubvariety.MultiplierBaseIdeals.coneIdeal_mono, EnvelopeProperty.ProjSubvariety.MultiplierBaseIdeals.parameterObstacleIdeal_map_cone, EnvelopeProperty.ProjSubvariety.MultiplierBaseIdeals.coneObstacleIdeal_le_coneIdeal_of_schemeIdeal_le}
\leanok
\uses{prop:c0n-original-obstacle-scheme-sequence}
For an actual whole ideal $K$ on the same family, push it along the original
closed inclusion, take the native homogeneous representative from all original
coordinate ideals, and map it by $k[t][x]\to (k[x]/I)[t]$. This defines a
canonical original cone ideal $\mathcal B(K)$. The homogeneous representative's
whole sheaf pulls back to exactly $K$, and $\mathcal B$ is monotone. If the
actual original $\mathfrak a_m$ is contained in $K$, then the same original
$J_m(\phi)$ is contained in $\mathcal B(K)$.
\end{lemma}
\begin{proof}
The native representative is the homogeneous core of the intersection of all
actual coordinate-ideal preimages. It is constructed, with no representation
witness. Its entire quotient-kernel sheaf is the original ambient ideal. The
actual closed-immersion ideal extension retracts on pullback, giving exactly
$K$. Both coordinate preimage and homogeneous core are monotone. For a
homogeneous ideal $H$, an inclusion of its whole sheaf in the ambient ideal
puts every dehomogenization of $H$ in every coordinate ideal. The native
homogeneous-core Galois connection therefore gives $H$ contained in that
representative. Apply the actual ideal pull/push adjunction to
$\mathfrak a_m\subseteq K$, then apply this result to the original
parameter-swapped $J_m$. Composing the genuine parameter exchange with its
inverse before the original cone quotient is exactly the original cone
quotient. Thus taking ideal images gives the asserted original inclusion.
\end{proof}

\begin{lemma}[Whole scheme multiplication yields original cone saturation]
\label{lem:c0n-scheme-cone-saturation-product}
\lean{EnvelopeProperty.ProjSubvariety.MultiplierBaseIdeals.familyIdeal_le_ker_parameterConeQuotient, EnvelopeProperty.ProjSubvariety.MultiplierBaseIdeals.parameterConeQuotient_X, EnvelopeProperty.ProjSubvariety.MultiplierBaseIdeals.mem_saturation_of_coordinate_powers, EnvelopeProperty.ProjSubvariety.MultiplierBaseIdeals.coordinate_power_mem_of_sheaf_le, EnvelopeProperty.ProjSubvariety.MultiplierBaseIdeals.ambient_sheaf_le_of_family_comap_le, EnvelopeProperty.ProjSubvariety.MultiplierBaseIdeals.map_cone_le_saturation_of_family_comap_le, EnvelopeProperty.ProjSubvariety.MultiplierBaseIdeals.homogeneousPower, EnvelopeProperty.ProjSubvariety.MultiplierBaseIdeals.homogeneousPower_toIdeal, EnvelopeProperty.ProjSubvariety.MultiplierBaseIdeals.quotientIdealSheaf_comap_pow, EnvelopeProperty.ProjSubvariety.MultiplierBaseIdeals.coneIdeal_pow_le_saturation}
\leanok
\uses{lem:c0n-canonical-scheme-cone-representative,lem:c0n-homogeneous-sheaf-arithmetic}
For whole ideals on the actual original family,
$\mathcal B(K_1K_2)\subseteq(\mathcal B(K_1)\mathcal B(K_2))^{\rm sat}$.
\end{lemma}
\begin{proof}
Let $H_1,H_2$ be their constructed original homogeneous representatives. Their
actual sheaves pull back to $K_1,K_2$, so the sheaf of $H_1H_2$ pulls back to
$K_1K_2$ by the proved whole-sheaf product law. A homogeneous numerator in the
representative of $K_1K_2$ is therefore locally in $H_1H_2$ modulo the original
family defining ideal. On every original coordinate chart the native quotient
away criterion supplies a power of that same coordinate taking the numerator
into $H_1H_2$ plus the original defining ideal. There are finitely many original
coordinates. A power of the ideal they generate is contained in the ideal of
these coordinate powers (expand its monomial generators and use the
pigeonhole principle). Hence a single irrelevant-ideal power times the
numerator lies in that sum. The original cone quotient kills the same defining
ideal and sends the coordinate ideal to the same original cone irrelevant
ideal. This is precisely the native definition of irrelevant saturation.
Decomposing an arbitrary element into its finitely many homogeneous
components and taking the largest required power proves the whole ideal
inclusion. This is an actual localization-and-saturation obligation, not a
comparison premise in the final envelope theorem.
\end{proof}

\begin{lemma}[Actual closed extension adds exactly the original kernel]
\label{lem:c0n-original-closed-extension-kernel}
\lean{EnvelopeProperty.Char0.EnvelopeGradedBaseIdealsClosedExtension.map_comap_eq_sup_ker, EnvelopeProperty.Char0.EnvelopeGradedBaseIdealsClosedExtension.quotientIdealSheaf_sup}
\leanok
For an actual closed immersion $i:Y\to P$ and a whole original ideal $I$ on
$P$, $(i^*I)_* = I+\ker i$. Homogeneous ideal sheafification preserves sums.
\end{lemma}
\begin{proof}
Every original affine open of $P$ has affine preimage, and the same original
section map is surjective. Native geometric ideal comap is ring-ideal image;
actual ideal extension is its ring-ideal preimage. The surjective ring ideal
correspondence gives preimage of image equal to the original ideal plus the
original section-map kernel. The native kernel ideal affine formula proves
the whole ideal equality. For homogeneous sums, exact dehomogenization on
the original coordinate cover and preservation of sums by ideal images give
the whole sheaf equality.
\end{proof}

\begin{proposition}[The actual original asymptotic multiplier cone sequence]
\label{prop:c0n-original-asymptotic-multiplier-cone}
\lean{EnvelopeProperty.ProjSubvariety.MultiplierBaseIdeals.familyAbsoluteStructure, EnvelopeProperty.ProjSubvariety.MultiplierBaseIdeals.schemeMultiplierSequence, EnvelopeProperty.ProjSubvariety.MultiplierBaseIdeals.schemeIdeal_le_schemeMultiplierSequence, EnvelopeProperty.ProjSubvariety.MultiplierBaseIdeals.schemeMultiplierSequence_mul_le_pow, EnvelopeProperty.ProjSubvariety.MultiplierBaseIdeals.multiplierConeIdeal, EnvelopeProperty.ProjSubvariety.MultiplierBaseIdeals.coneObstacleIdeal_le_multiplierConeIdeal, EnvelopeProperty.ProjSubvariety.MultiplierBaseIdeals.multiplierConeIdeal_mul_le_saturation}
\leanok
\uses{prop:c0n-original-obstacle-generic-stalk,lem:c0n-canonical-scheme-cone-representative,lem:c0n-scheme-cone-saturation-product}
On the same original smooth integral parameter family over the same
characteristic-zero field, form the actual whole asymptotic multiplier ideals
of the original $\mathfrak a_\bullet$, and define $B_m$ by the preceding
canonical cone construction. Then $J_m(\phi)\subseteq B_m$ for $m>0$, and
$B_{ml}\subseteq(B_m^l)^{\rm sat}$ for all $m,l$, including zero indices.
No vanishing, generation, target ideal comparison, or saturation statement
is an input.
\end{proposition}
\begin{proof}
The already produced original generic nonvanishing defines the full actual
scheme asymptotic multiplier ideal without a nonvanishing premise. Native
whole scheme multiplier containment puts $\mathfrak a_m$ in this same ideal,
which yields the original cone containment by the canonical construction.
Native whole scheme asymptotic subadditivity puts the index $ml$ ideal in the
$l$th power of the same index $m$ ideal. Monotonicity of the actual cone
construction followed by the proved whole-scheme-power saturation comparison
gives the required bound. Homogeneous powers are constructed recursively by
the actual homogeneous product and unit; their underlying ideals are exactly
the original ordinary ideal powers. This includes $l=0$ and requires no
unproved homogeneous-power instance.
\end{proof}

\begin{corollary}[Original algebraic smooth-chart scope]
\label{cor:c0n-original-multiplier-cone-smooth-charts}
\lean{EnvelopeProperty.ProjSubvariety.MultiplierBaseIdeals.originalMultiplierConeIdeal, EnvelopeProperty.ProjSubvariety.MultiplierBaseIdeals.coneObstacleIdeal_le_originalMultiplierConeIdeal, EnvelopeProperty.ProjSubvariety.MultiplierBaseIdeals.originalMultiplierConeIdeal_mul_le_saturation}
\leanok
\uses{prop:c0n-original-asymptotic-multiplier-cone}
For the original $X$ with its original smooth-chart condition and original
lower-bounded obstacle, the same actually constructed $B_m$ satisfies the
preceding containments. No smooth-family, integral-family, dimension,
nonvanishing, generation, or ideal-comparison certificate is supplied.
\end{corollary}
\begin{proof}
The native smoothness producer constructs the actual original family product
isomorphism and proves smoothness of its literal original total field
structure from the original algebraic smooth charts. Install that proved
smoothness for this same field structure and use the preceding construction.
The family integrality and every original generic stalk were already produced
from the original equations and obstacle. Thus all extra geometric plumbing
is internal and the original cone and ideals are unchanged.
\end{proof}

\begin{lemma}[The actual parameter coordinate quotient]
\label{lem:c0n-original-parameter-chart-quotient}
\lean{EnvelopeProperty.ProjSubvariety.MultiplierBaseIdeals.chartParameterEquiv, EnvelopeProperty.ProjSubvariety.MultiplierBaseIdeals.chartParameterEquiv_map_C, EnvelopeProperty.ProjSubvariety.MultiplierBaseIdeals.familyCoordinateIdeal_map, EnvelopeProperty.ProjSubvariety.MultiplierBaseIdeals.quotientPolynomialAlgEquiv, EnvelopeProperty.ProjSubvariety.MultiplierBaseIdeals.familyCoordinateQuotientEquiv}
\leanok
The whole native affine coordinate quotient of the original parameter family
is canonically the polynomial ring over the original quotient chart $A_j$,
as an algebra over the original coefficient field. Its constructed maps use
the actual coefficient-variable interchange and original quotient maps.
\end{lemma}
\begin{proof}
On the complementary coordinate variables, use the actual native
coefficient-variable interchange $k[t][z]\simeq k[z][t]$. Coefficient extension
and original dehomogenization commute, so the complete original family
coordinate ideal maps exactly to the polynomial extension of the complete
original coordinate ideal $I_j$. The genuine quotient-ring isomorphism then
gives $k[z][t]/I_jk[z][t]\simeq(k[z]/I_j)[t]$. Finally apply the already proved
actual original coordinate quotient equivalence $k[z]/I_j\simeq A_j$
coefficientwise to the whole polynomial ring. The native quotient maps and
polynomial map identities prove that every original homogeneous numerator,
coefficient, and parameter has its same original image. No ideal or chart
comparison is supplied as a certificate.
\end{proof}

\begin{lemma}[Actual local section algebras and original scalar triangles]
\label{lem:c0n-original-section-algebra-coherence}
\lean{EnvelopeProperty.Char0.EnvelopeGradedBaseIdealsSectionAlgebra.sectionBaseMap_eq_sectionsAlgebra, EnvelopeProperty.Char0.EnvelopeGradedBaseIdealsSectionAlgebra.sectionBaseMap_comp, EnvelopeProperty.Char0.EnvelopeGradedBaseIdealsSectionAlgebra.sectionAlgEquivOfIso}
\leanok
The native section-base map is the same original map defining the regular-smooth
section algebra. Composing the original structure morphism with a coefficient
spectrum map composes these same section-base ring maps. An actual scheme
isomorphism satisfying its true original scalar triangle induces an algebra
isomorphism on the corresponding whole local section rings.
\end{lemma}
\begin{proof}
Expand the two named section maps: both are the original inverse spectrum
section isomorphism followed by the original structure pullback and original
restriction. For composition, native spectrum-section naturality moves the
original coefficient homomorphism through this same inverse isomorphism.
The structure-sheaf composition law and associativity give the entire ring-map
equality. For a scheme isomorphism its whole section map is an isomorphism;
the original scalar triangle and native section-base naturality prove that
this same ring isomorphism is linear. No substitute scalar action is used.
\end{proof}

\begin{lemma}[The same whole original family affine chart]
\label{lem:c0n-original-family-affine-sections}
\lean{EnvelopeProperty.ProjSubvariety.MultiplierBaseIdeals.familyCutIdeal, EnvelopeProperty.ProjSubvariety.MultiplierBaseIdeals.familyCutIso, EnvelopeProperty.ProjSubvariety.MultiplierBaseIdeals.familyCutIso_inclusion, EnvelopeProperty.ProjSubvariety.MultiplierBaseIdeals.familyCutProjection, EnvelopeProperty.ProjSubvariety.MultiplierBaseIdeals.familyCutIso_projection, EnvelopeProperty.ProjSubvariety.MultiplierBaseIdeals.familyCutChartOpen, EnvelopeProperty.ProjSubvariety.MultiplierBaseIdeals.familyChartOpen, EnvelopeProperty.ProjSubvariety.MultiplierBaseIdeals.familyChartOpen_eq, EnvelopeProperty.ProjSubvariety.MultiplierBaseIdeals.familyChartOpen_isAffine, EnvelopeProperty.ProjSubvariety.MultiplierBaseIdeals.familyChartCoordinateAlgEquiv, EnvelopeProperty.ProjSubvariety.MultiplierBaseIdeals.familyChartSectionsEquiv, EnvelopeProperty.ProjSubvariety.MultiplierBaseIdeals.familyChartOpen_nonempty, EnvelopeProperty.ProjSubvariety.MultiplierBaseIdeals.familyChartAffineOpen}
\leanok
\uses{lem:c0n-original-parameter-chart-quotient,lem:c0n-original-section-algebra-coherence,lem:c0n-original-family-integral}
For every original relevant coordinate, the same actual family inverse
coordinate open is affine and nonempty, and its whole section ring is
canonically $A_j[t]$ as an algebra for the SAME actual total field structure.
\end{lemma}
\begin{proof}
The original family closed inclusion is isomorphic to its actual image
subscheme. Its image ideal is precisely the same original homogeneous
quotient kernel. Pull back its genuine coordinate affine open along this
actual isomorphism. The original inclusion triangle identifies the result
with the original family coordinate inverse image. Closed-immersion
preimages and isomorphism preimages preserve affineness. The actual zero
homogeneous prime belongs to this open because the original coordinate
stays outside the original family defining ideal, proving nonemptiness.
On the image subscheme the native whole coordinate section equivalence gives
the full original dehomogenized quotient. The actual scheme isomorphism
transports whole sections, with its proved original parameter projection
triangle. Compose with the preceding true polynomial-quotient equivalence.
Its parameter-linear source action restricts to the original total-field
action: the actual total structure is the original parameter structure
followed by the original coefficient spectrum map, so section-base
composition proves the scalar tower. This is exactly the section algebra
used in the original scheme multiplier/DVR theorem.
\end{proof}

\begin{lemma}[All original parameter-polynomial coordinate values]
\label{lem:c0n-original-parameter-coordinate-values}
\lean{EnvelopeProperty.ProjSubvariety.MultiplierBaseIdeals.quotientPolynomialAlgEquiv_mk, EnvelopeProperty.ProjSubvariety.MultiplierBaseIdeals.familyCoordinateQuotientEquiv_mk, EnvelopeProperty.ProjSubvariety.MultiplierBaseIdeals.coordinateQuotientEquiv_mk_dehom, EnvelopeProperty.ProjSubvariety.MultiplierBaseIdeals.familyCoordinateQuotientEquiv_parameterDehom}
\leanok
\uses{lem:c0n-original-parameter-chart-quotient}
The constructed family coordinate quotient equivalence sends the class of
the dehomogenization of the parameter exchange of EVERY original polynomial
$g\in k[x][t]$ to the SAME original coefficientwise dehomogenization
$g(\delta_jx,t)\in A_j[t]$.
\end{lemma}
\begin{proof}
The native polynomial quotient equivalence sends a quotient class to the
coefficientwise quotient polynomial. The original coordinate quotient
isomorphism sends the dehomogenization of every original variable to the
same old quotient-chart ratio, and fixes the original field. Algebra-map
extensionality therefore gives this equality on every original coefficient
polynomial. For $s t^l$, the actual parameter exchange is $t^l s$;
coefficient interchange after dehomogenization sends this to
$\delta_js\,t^l$. The quotient polynomial and original coordinate equivalences
act coefficientwise, so the entire composite is the same original
$\delta_js\,t^l$. Additivity and the original finite polynomial expansion
prove the assertion for all $g$, including zero and inhomogeneous elements.
\end{proof}

\begin{lemma}[The entire original chart coefficient diagram]
\label{lem:c0n-original-native-chart-coefficients}
\lean{EnvelopeProperty.ProjSubvariety.MultiplierBaseIdeals.familyCoordinateSection, EnvelopeProperty.ProjSubvariety.MultiplierBaseIdeals.familyCoordinateSection_eq_transport, EnvelopeProperty.ProjSubvariety.MultiplierBaseIdeals.familyChartCoordinateAlgEquiv_coordinateSection, EnvelopeProperty.ProjSubvariety.MultiplierBaseIdeals.familyChartSectionsEquiv_coordinateSection, EnvelopeProperty.ProjSubvariety.MultiplierBaseIdeals.familyChartSectionsEquiv_parameterDehom, EnvelopeProperty.ProjSubvariety.MultiplierBaseIdeals.familyCoordinateSection_mem_of_representative, EnvelopeProperty.ProjSubvariety.MultiplierBaseIdeals.familyChartSectionsEquiv_symm_dehom_mem, EnvelopeProperty.ProjSubvariety.MultiplierBaseIdeals.schemeIdeal_chart_map}
\leanok
\uses{lem:c0n-original-family-affine-sections,lem:c0n-original-parameter-coordinate-values}
The actual family section equivalence sends the original pulled coordinate
section of the dehomogenization of $E(g)$ to the same original
$g(\delta_jx,t)$. If $P$ lies in the canonical representative of a whole family
ideal $K$, the inverse image of the same original
$E^{-1}(P)(\delta_jx,t)$ belongs to $K$ on this actual affine chart.
The chart ideal of the original $\mathfrak a_q$, mapped by this same section
equivalence, is exactly the original obstacle ideal mapped by coefficientwise
$\delta_j$.
\end{lemma}
\begin{proof}
Take the actual original inclusion pullback on the fixed genuine family
coordinate open, expressed with the native appLE map. The proved actual
inclusion triangle identifies it with the composite of the original image
inclusion section map and the actual family-to-image section isomorphism.
The native quotient section equivalence sends that original composite to the
original quotient class. The preceding parameter-polynomial formula gives
the claimed image for every original $g$. For membership, the canonical
representative places the dehomogenization in the original ambient coordinate
ideal of $K_*$. The actual affine ideal-comap formula sends it into the
pullback of $K_*$, which is exactly $K$ by the genuine closed-immersion
retraction. Apply the inverse of the same section equivalence. Finally the
whole original $\mathfrak a_q$ is the actual comap of the original parameter
obstacle ideal sheaf. Its exact native coordinate ideal is the full
parameter obstacle ideal mapped by dehomogenization and the original
coordinate section isomorphism. Compose all actual ring maps and use the
proved all-polynomial identity; ring-ideal image composition proves the
entire asserted ideal equality, with no coefficient-subfamily replacement.
\end{proof}
